\documentclass[10pt,a4paper]{amsart}
\usepackage{amsmath,amssymb,amsthm,graphicx}
% Retained minimal preamble for the Stacks Project chapter more-morphisms.tex (More on Morphisms).
% Verbatim excerpt of the chapter preamble: Xy-pic setup (native preamble lines 14-18) and the native
% theorem declarations with the mathematical macros (native preamble lines 154-217). Copied mechanically;
% nothing added, removed or redefined. The wrapper class and the amsmath/amssymb/amsthm/graphicx base
% packages are supplied by the assembler, so the class/externaldocument part of the native preamble is omitted.
% For commutative diagrams we use Xy-pic
\usepackage[all]{xy}

% We use 2cell for 2-commutative diagrams.
\xyoption{2cell}
% Theorem environments.
%
\theoremstyle{plain}
\newtheorem{theorem}[subsection]{Theorem}
\newtheorem{proposition}[subsection]{Proposition}
\newtheorem{lemma}[subsection]{Lemma}

\theoremstyle{definition}
\newtheorem{definition}[subsection]{Definition}
\newtheorem{example}[subsection]{Example}
\newtheorem{exercise}[subsection]{Exercise}
\newtheorem{situation}[subsection]{Situation}

\theoremstyle{remark}
\newtheorem{remark}[subsection]{Remark}
\newtheorem{remarks}[subsection]{Remarks}

\numberwithin{equation}{subsection}

% Macros
%
\def\lim{\mathop{\mathrm{lim}}\nolimits}
\def\colim{\mathop{\mathrm{colim}}\nolimits}
\def\Spec{\mathop{\mathrm{Spec}}}
\def\Hom{\mathop{\mathrm{Hom}}\nolimits}
\def\Ext{\mathop{\mathrm{Ext}}\nolimits}
\def\SheafHom{\mathop{\mathcal{H}\!\mathit{om}}\nolimits}
\def\SheafExt{\mathop{\mathcal{E}\!\mathit{xt}}\nolimits}
\def\Sch{\mathit{Sch}}
\def\Mor{\mathop{\mathrm{Mor}}\nolimits}
\def\Ob{\mathop{\mathrm{Ob}}\nolimits}
\def\Sh{\mathop{\mathit{Sh}}\nolimits}
\def\NL{\mathop{N\!L}\nolimits}
\def\CH{\mathop{\mathrm{CH}}\nolimits}
\def\proetale{{pro\text{-}\acute{e}tale}}
\def\etale{{\acute{e}tale}}
\def\QCoh{\mathit{QCoh}}
\def\Ker{\mathop{\mathrm{Ker}}}
\def\Im{\mathop{\mathrm{Im}}}
\def\Coker{\mathop{\mathrm{Coker}}}
\def\Coim{\mathop{\mathrm{Coim}}}

% Boxtimes
%
\DeclareMathSymbol{\boxtimes}{\mathbin}{AMSa}{"02}

%
% Macros for moduli stacks/spaces
%
\def\QCohstack{\mathcal{QC}\!\mathit{oh}}
\def\Cohstack{\mathcal{C}\!\mathit{oh}}
\def\Spacesstack{\mathcal{S}\!\mathit{paces}}
\def\Quotfunctor{\mathrm{Quot}}
\def\Hilbfunctor{\mathrm{Hilb}}
\def\Curvesstack{\mathcal{C}\!\mathit{urves}}
\def\Polarizedstack{\mathcal{P}\!\mathit{olarized}}
\def\Complexesstack{\mathcal{C}\!\mathit{omplexes}}
% \Pic is the operator that assigns to X its picard group, usage \Pic(X)
% \Picardstack_{X/B} denotes the Picard stack of X over B
% \Picardfunctor_{X/B} denotes the Picard functor of X over B
\def\Pic{\mathop{\mathrm{Pic}}\nolimits}
\def\Picardstack{\mathcal{P}\!\mathit{ic}}
\def\Picardfunctor{\mathrm{Pic}}
\def\Deformationcategory{\mathcal{D}\!\mathit{ef}}
\newenvironment{reference}{}{}
\newenvironment{slogan}{}{}
\newenvironment{history}{}{}
\usepackage{cleveref}
\crefname{originallemma}{original lemma}{original lemmas}
\Crefname{originallemma}{Original lemma}{Original lemmas}
\makeatletter
\newlabel{algebra-lemma-NL-localize-bottom}{{07BS}{source}{}{sourcepointer-1}{}}
\newlabel{algebra-lemma-NL-localize-bottom@cref}{{[originallemma][0][]07BS}{[1][0][]source}{}{}{}}
\makeatother
\crefname{originallemma}{original lemma}{original lemmas}
\Crefname{originallemma}{Original lemma}{Original lemmas}
\makeatletter
\newlabel{algebra-lemma-localize-NL}{{00S7}{source}{}{sourcepointer-2}{}}
\newlabel{algebra-lemma-localize-NL@cref}{{[originallemma][0][]00S7}{[1][0][]source}{}{}{}}
\makeatother
\crefname{originallemma}{original lemma}{original lemmas}
\Crefname{originallemma}{Original lemma}{Original lemmas}
\makeatletter
\newlabel{modules-lemma-stalk-NL}{{0D09}{source}{}{sourcepointer-3}{}}
\newlabel{modules-lemma-stalk-NL@cref}{{[originallemma][0][]0D09}{[1][0][]source}{}{}{}}
\makeatother
\crefname{originallemma}{original lemma}{original lemmas}
\Crefname{originallemma}{Original lemma}{Original lemmas}
\makeatletter
\newlabel{modules-section-netherlander}{{08TG}{source}{}{sourcepointer-4}{}}
\newlabel{modules-section-netherlander@cref}{{[originallemma][0][]08TG}{[1][0][]source}{}{}{}}
\makeatother
\crefname{originallemma}{original lemma}{original lemmas}
\Crefname{originallemma}{Original lemma}{Original lemmas}
\makeatletter
\newlabel{schemes-lemma-compare-constructions}{{01I7}{source}{}{sourcepointer-5}{}}
\newlabel{schemes-lemma-compare-constructions@cref}{{[originallemma][0][]01I7}{[1][0][]source}{}{}{}}
\makeatother
\crefname{originallemma}{original lemma}{original lemmas}
\Crefname{originallemma}{Original lemma}{Original lemmas}
\makeatletter
\newlabel{schemes-lemma-spec-sheaves}{{01HV}{source}{}{sourcepointer-6}{}}
\newlabel{schemes-lemma-spec-sheaves@cref}{{[originallemma][0][]01HV}{[1][0][]source}{}{}{}}
\makeatother
\crefname{originallemma}{original lemma}{original lemmas}
\Crefname{originallemma}{Original lemma}{Original lemmas}
\makeatletter
\newlabel{sheaves-lemma-stalk-pullback}{{008H}{source}{}{sourcepointer-7}{}}
\newlabel{sheaves-lemma-stalk-pullback@cref}{{[originallemma][0][]008H}{[1][0][]source}{}{}{}}
\makeatother
\title{The naive cotangent complex of an affine chart computes that of the morphism}
\author{The Stacks Project Authors}
\date{stacks-project:more-morphisms.tex:local-snapshot-20261011; GNU FDL 1.2 or later}
\begin{document}
\maketitle
\noindent\textit{Editorial note.} Statement span is lines 3735-3745 of sources/stacks-more-morphisms/pinned-original-source.tex; proof span is lines 3747-3774, the complete author proof (26 body lines).\par\medskip
\noindent\textit{Editorial note.} Navigation stubs: 7 cited-but-not-included labels are supplied as non-load-bearing pointers with their pinned tags where the snapshot tag index has them.\par\medskip
\noindent\textit{Editorial note.} extra\_packages is empty on purpose: the retained preamble already loads xy with options and a second declaration clashes in the assembler wrapper.\par\medskip
\noindent\textit{Editorial note.} Text is retained verbatim, including the author's own typos and placeholders.\par\medskip
\noindent\textit{Allowed standard prerequisites (editorial).}\begin{itemize}
\item Explicit prerequisite boundary. The naive cotangent complex of a ring map and of a morphism, presentations, and sheafification; the comparison is retained.
\end{itemize}\medskip
\ifdefined\hypertarget\hypertarget{sourcepointer-1}{}\fi
\noindent\textit{Source reference.} Original lemma 07BS: sources/stacks-more-morphisms/pinned-original-source.tex, source lines 3774--3774. Cited by the retained statement or proof but not part of the collected spans; supplied only as a navigation pointer. The argument does not depend on its content.\par\medskip
\ifdefined\hypertarget\hypertarget{sourcepointer-2}{}\fi
\noindent\textit{Source reference.} Original lemma 00S7: sources/stacks-more-morphisms/pinned-original-source.tex, source lines 3774--3774. Cited by the retained statement or proof but not part of the collected spans; supplied only as a navigation pointer. The argument does not depend on its content.\par\medskip
\ifdefined\hypertarget\hypertarget{sourcepointer-3}{}\fi
\noindent\textit{Source reference.} Original lemma 0D09: sources/stacks-more-morphisms/pinned-original-source.tex, source lines 3774--3774. Cited by the retained statement or proof but not part of the collected spans; supplied only as a navigation pointer. The argument does not depend on its content.\par\medskip
\ifdefined\hypertarget\hypertarget{sourcepointer-4}{}\fi
\noindent\textit{Source reference.} Original lemma 08TG: sources/stacks-more-morphisms/pinned-original-source.tex, source lines 3774--3774. Cited by the retained statement or proof but not part of the collected spans; supplied only as a navigation pointer. The argument does not depend on its content.\par\medskip
\ifdefined\hypertarget\hypertarget{sourcepointer-5}{}\fi
\noindent\textit{Source reference.} Original lemma 01I7: sources/stacks-more-morphisms/pinned-original-source.tex, source lines 3774--3774. Cited by the retained statement or proof but not part of the collected spans; supplied only as a navigation pointer. The argument does not depend on its content.\par\medskip
\ifdefined\hypertarget\hypertarget{sourcepointer-6}{}\fi
\noindent\textit{Source reference.} Original lemma 01HV: sources/stacks-more-morphisms/pinned-original-source.tex, source lines 3774--3774. Cited by the retained statement or proof but not part of the collected spans; supplied only as a navigation pointer. The argument does not depend on its content.\par\medskip
\ifdefined\hypertarget\hypertarget{sourcepointer-7}{}\fi
\noindent\textit{Source reference.} Original lemma 008H: sources/stacks-more-morphisms/pinned-original-source.tex, source lines 3774--3774. Cited by the retained statement or proof but not part of the collected spans; supplied only as a navigation pointer. The argument does not depend on its content.\par\medskip
\begin{lemma}
\label{lemma-NL-affine}
Let $f : X \to Y$ be a morphism of schemes. Let
$\Spec(A) = U \subset X$ and $\Spec(R) = V \subset S$
be affine opens with $f(U) \subset V$.
There is a canonical map
$$
\widetilde{\NL_{A/R}} \longrightarrow \NL_{X/Y}|_U
$$
of complexes which is an isomorphism in $D(\mathcal{O}_U)$.
\end{lemma}

\begin{proof}
From the construction of $\NL_{X/Y}$ in
Modules, Section \ref{modules-section-netherlander}
we see there is a canonical map of complexes
$\NL_{\mathcal{O}_X(U)/f^{-1}\mathcal{O}_Y(U)} \to \NL_{X/Y}(U)$
of $A = \mathcal{O}_X(U)$-modules, which is compatible
with further restrictions. Using the canonical map
$R \to f^{-1}\mathcal{O}_Y(U)$ we obtain a canonical map
$\NL_{A/R} \to \NL_{\mathcal{O}_X(U)/f^{-1}\mathcal{O}_Y(U)}$
of complexes of $A$-modules.
Using the universal property of the $\widetilde{\ }$
functor (see Schemes, Lemma \ref{schemes-lemma-compare-constructions})
we obtain a map as in the statement of the lemma.
We may check this map is an isomorphism on cohomology sheaves
by checking it induces isomorphisms on stalks.
This follows from
Algebra, Lemma \ref{algebra-lemma-NL-localize-bottom} and
\ref{algebra-lemma-localize-NL}
and
Modules, Lemma \ref{modules-lemma-stalk-NL}
(and the description of the stalks of
$\mathcal{O}_X$ and $f^{-1}\mathcal{O}_Y$
at a point $\mathfrak p \in \Spec(A)$ as $A_\mathfrak p$ and
$R_\mathfrak q$ where $\mathfrak q = R \cap \mathfrak p$; references
used are Schemes, Lemma \ref{schemes-lemma-spec-sheaves}
and
Sheaves, Lemma \ref{sheaves-lemma-stalk-pullback}).
\end{proof}

\end{document}
